\chapter{New Revolutions in Particle Physics: The Standard Model}

This lecture begins with a joke, but the mathematics begins with a sharp question: what is the relation between an abstract symmetry group and the concrete matrices that act on physical states? We answer that question first for rotations and spin, then repeat the same pattern for color \(SU(3)\), and only after that turn to the dynamics that single out colorless states as the particles we actually observe. These notes follow Leonard Susskind's lecture closely and were curated by LazyingArt LLC.

\section{From Abstract Groups to Physical Representations}

The lecture re-enters the main storyline through a basic but important distinction. There is, on the one hand, the abstract notion of a symmetry operation. A rotation in space is a perfect example: as an abstract group element it is simply an operation that can be performed. On the other hand, physics asks how that abstract operation acts on the variables or states that matter physically.

For classical mechanics, this distinction is easy to see. A rotation acts on coordinates, or on the direction of a magnet, or on any other spatial object. In quantum mechanics the point becomes sharper, because the states of a system are represented by vectors in a linear vector space. The symmetry is abstract; the representation is the matrix action on that vector space.

\subsection{Question \& Answer}

\paragraph{Question.} Can we dissociate the abstract group from the matrices that act on states?

\paragraph{Answer.} We should distinguish them, but not separate them. The abstract group is the symmetry itself, with its multiplication law. A representation is the way that symmetry acts on a particular system. The same abstract rotation may be represented by different matrices on different state spaces, but the matrices must reproduce the same group multiplication. In that sense the abstract group and the representation are different notions, yet in physics they are tied together from the start.

\section{Why Representation Dimension Depends on the System}

Once we think in terms of states, the size of the matrices is no longer mysterious. It is determined by the number of mutually orthogonal states available to the system under discussion.

For spin alone this is familiar:
\begin{equation}
\text{spin }0:\ 1\times1,\qquad
\text{spin }\frac12:\ 2\times2,\qquad
\text{spin }1:\ 3\times3,\qquad
\text{spin }\frac32:\ 4\times4.
\end{equation}
A spin-zero system is trivial under rotations: the matrix is just the unit \(1\times1\) matrix. A spin-\(\frac12\) system has two orthogonal states, spin up and spin down, so rotations act by \(2\times2\) matrices. A spin-one system has three orthogonal states, so rotations act by \(3\times3\) matrices.

The same abstract rotation may therefore have many different matrix realizations. That is not a side remark. It is the main lesson that prepares the passage from \(SU(2)\), which we know from spin, to \(SU(3)\), which we will use for color.

The lecture also stresses a more extreme example. If we consider not just the spin of an electron but its orbital motion, then the number of mutually orthogonal states is infinite. The state vector may be described by amplitudes for every possible position, or for every possible momentum, and therefore the corresponding representation of the rotation group is infinite-dimensional. The abstract group has not changed at all; only the system on which it acts has changed.

\section{Special Unitary Matrices and Infinitesimal Generators}

Let us now return to the concrete case of a spin-\(\frac12\) particle. Rotations are represented by \(2\times2\) unitary matrices \(U\). The unitary condition is
\begin{equation}
U^\dagger U = 1.
\end{equation}
A general \(2\times2\) complex matrix has four complex entries, hence eight real parameters. Unitarity imposes four real constraints, leaving four parameters. But a rotation in ordinary space depends only on three parameters. Therefore one more condition is needed:
\begin{equation}
\det U = 1.
\end{equation}
This is the special-unitary restriction. In compact form,
\begin{equation}
4\ \text{complex entries} = 8\ \text{real parameters}
\longrightarrow
4\ \text{after }U^\dagger U=1
\longrightarrow
3\ \text{after }\det U=1.
\end{equation}

The lecture pauses over the determinant because it is a consistent condition under multiplication. If \(A\) and \(B\) are matrices, then
\begin{equation}
\det(AB)=\det A\,\det B,
\end{equation}
and if \(U\) is unitary then \(\det U\) is a pure phase:
\begin{equation}
\det U = e^{i\theta}.
\end{equation}
Multiplying by an overall phase removes this extra parameter:
\begin{equation}
\det\!\bigl(e^{-i\theta/2}U\bigr)=1.
\end{equation}
This is why the rotation discussion may be reduced to special unitary matrices. The lecture then adds the usual caveat: this correspondence is not exactly one-to-one. The matrices \(U\) and \(-U\) correspond to the same spatial rotation, so \(SU(2)\) covers the ordinary rotation group twice.

Now the lecture shifts from finite group elements to infinitesimal ones. This is the central mathematical move, because the generators are the conserved quantities. We write a near-identity group element as
\begin{equation}
U = 1 + i\epsilon T,
\end{equation}
with \(\epsilon\) small.

\begin{figure}[t]
\centering
\includegraphics[width=0.58\textwidth]{lecture_04_figure_03.png}
\caption{First-order unitarity implies a Hermitian generator.}
\end{figure}

The board calculation is the one we want to preserve:
\begin{align}
(1+i\epsilon T)(1-i\epsilon T^\dagger) &= 1,\\
i\epsilon\,(T-T^\dagger) &= 0,\\
T &= T^\dagger.
\end{align}
So the infinitesimal generator is Hermitian. This is exactly what we should expect in quantum mechanics: Hermitian operators are observables.

At this stage the lecture briefly asks why there are three such generators for rotations. The answer is that infinitesimal rotations about the \(x\)-, \(y\)-, and \(z\)-axes form three linearly independent directions in the space of small rotations. More abstractly, the generators behave like the components of a vector, and in quantum mechanics they are the angular momentum operators.

But one condition has not yet been used. Since \(\det U = 1\), the infinitesimal generator must also be traceless. The lecture states the near-identity determinant expansion, and the board catches the moment at which it is being set up, even though the trace formula is not fully finished in the visible frame.

\begin{figure}[t]
\centering
\includegraphics[width=0.58\textwidth]{lecture_04_figure_02.png}
\caption{Near-identity determinant expansion on the board.}
\end{figure}

For a matrix close to the identity we use the standard first-order expansion
\begin{equation}
\det(1+\epsilon m)=1+\epsilon\,\operatorname{tr}m+O(\epsilon^2).
\end{equation}
Applying this to \(U=1+i\epsilon T\) gives
\begin{equation}
\det U = 1 + i\epsilon\,\operatorname{tr}T + O(\epsilon^2).
\end{equation}
Since \(\det U=1\), the first-order term must vanish:
\begin{equation}
\operatorname{tr}T=0.
\end{equation}

We have therefore arrived at the key algebraic statement for \(SU(2)\): the generators are precisely the traceless Hermitian \(2\times2\) matrices. There are three of them, and they are spanned by the Pauli matrices. In the language of the lecture, angular momentum in quantum mechanics is nothing but the generator of the rotation group.

\section{From \(SU(2)\) to \(SU(3)\)}

At this point the lecture says, in effect, that we now do the same mathematics again, but with \(3\times3\) matrices. A general \(3\times3\) complex matrix has nine complex entries, hence eighteen real parameters. Unitarity imposes nine real constraints, and the determinant-one condition removes one more. Thus
\begin{equation}
18 - 9 - 1 = 8.
\end{equation}
So \(SU(3)\) has eight independent generators, equivalently eight linearly independent traceless Hermitian \(3\times3\) matrices.

The physical interpretation is new even though the algebra is parallel. A quark carries three color amplitudes, and we may represent its color state by a column vector
\begin{equation}
q=
\begin{pmatrix}
q_r\\
q_b\\
q_g
\end{pmatrix}.
\end{equation}
A color rotation acts by
\begin{equation}
q' = Uq,
\end{equation}
where now \(U\in SU(3)\). In words, the matrix mixes red, blue, and green into quantum superpositions of one another.

Anti-quarks furnish the complex-conjugate representation. If \(q\) transforms by \(U\), then the anti-quark transforms by \(U^*\):
\begin{equation}
\bar q' = U^* \bar q.
\end{equation}
This is the \(\bar 3\), the anti-triplet. The lecture emphasizes that the \(3\) and the \(\bar 3\) are distinct representations of \(SU(3)\), unlike the special \(SU(2)\) case where the complex-conjugate representation is equivalent to the original one. It also states, at the lecture level, that the generators in the \(\bar 3\) representation are the negatives of those in the \(3\) representation, which is why the anti-quark carries the opposite color charge.

The eight traceless Hermitian \(3\times3\) generators are named as the Gell-Mann matrices. The lecture does not develop their explicit forms here. What matters is the count and the interpretation: eight generators, eight gluons, and one octet representation in which those generators transform among themselves.

\section{How Representations Combine}

Before turning to color, the lecture gives a warm-up from \(SU(2)\). If we combine two spin-\(\frac12\) systems, we obtain four states. But those four states do not remain an irreducible four-dimensional multiplet. They split into a singlet and a triplet:
\begin{equation}
2\otimes2 = 1\oplus3.
\end{equation}
In an explicit basis, and allowing ourselves the usual optional normalization factors, we may write
\begin{align}
\text{singlet:}\qquad
&\frac{1}{\sqrt2}\bigl(\psi_\uparrow\phi_\downarrow-\psi_\downarrow\phi_\uparrow\bigr),\\
\text{triplet:}\qquad
&\psi_\uparrow\phi_\uparrow,\quad
\frac{1}{\sqrt2}\bigl(\psi_\uparrow\phi_\downarrow+\psi_\downarrow\phi_\uparrow\bigr),\quad
\psi_\downarrow\phi_\downarrow.
\end{align}
The important structural fact is that the antisymmetric combination does not mix with the other three under the group action. It forms a one-dimensional invariant subspace. The remaining three states mix among themselves and form the spin-one representation.

\subsection{Question \& Answer}

\paragraph{Question.} What does an equation such as \(3\otimes3=\bar 3\oplus6\) actually mean?

\paragraph{Answer.} It means that the composite state space built from two triplets is not irreducible. The tensor-product space has dimension \(3\times3=9\), but in a suitable basis the group action breaks into smaller invariant blocks. The symbol \(\oplus\) means that the representation decomposes into irreducible pieces. The lecture even says, in plain language, that the plus sign means \emph{or}: the two-particle system transforms either as one irreducible multiplet or as the other, depending on which linear combination of states we choose.

This is why the late question about block diagonal form is so useful. In the right basis, a large representation matrix separates into blocks, one for each irreducible sector. The singlet does not mix with the triplet in the \(2\otimes2\) example, and exactly the same logic governs the color products.

We now apply that logic to \(SU(3)\).

\begin{figure}[t]
\centering
\includegraphics[width=0.62\textwidth]{lecture_04_figure_04.png}
\caption{Two-quark triplet product decomposed into \(\bar 3\) and \(6\).}
\end{figure}

For two quarks,
\begin{equation}
3\otimes3=\bar 3\oplus6.
\end{equation}
The lecture interprets this in the simplest possible way. If we think of the quarks as carrying color indices, then the two-quark states may be combined symmetrically or antisymmetrically. The antisymmetric combinations furnish the \(\bar 3\); the symmetric combinations furnish the \(6\).

For a quark and an anti-quark, the lecture gives
\begin{equation}
3\otimes\bar 3 = 1\oplus8.
\end{equation}
The singlet is the colorless combination, schematically
\begin{equation}
q_r\bar q_r + q_b\bar q_b + q_g\bar q_g,
\end{equation}
up to normalization. The remaining eight states form the octet, which is the adjoint representation: they transform in the same way as the eight generators.

For three quarks, the lecture proceeds step by step. First combine two quarks:
\begin{equation}
3\otimes3=\bar 3\oplus6.
\end{equation}
Then combine with the third quark. Since
\begin{equation}
\bar 3\otimes3 = 1\oplus8,
\end{equation}
the three-quark product contains a singlet. The lecture writes this singlet with the antisymmetric epsilon tensor:
\begin{equation}
\epsilon_{ijk}q_i q_j q_k.
\end{equation}
Because \(\epsilon_{ijk}\) vanishes unless all three indices are different, this is precisely the antisymmetrized red--green--blue combination.

\section{Singlets, Hadrons, and the Octet of Gluons}

Once the representation theory has been laid out, the lecture translates it into particle language. A color-singlet \(qqq\) state is a baryon. A color-singlet \(q\bar q\) state is a meson. A gluon transforms in the octet representation, not in the singlet, because its color quantum numbers are those of a quark and an anti-quark with the singlet piece removed.

This is why the lecture insists on the octet of gluons. The quark transforms as a \(3\), the anti-quark as a \(\bar 3\), and the gluon field as an \(8\). The omitted singlet piece would not mix with the rest under \(SU(3)\), so dropping it is mathematically consistent.

The same logic leads naturally to glueballs. Since a single gluon is not a color singlet, it cannot appear as a free particle. But two gluons may combine into a singlet.

\begin{figure}[t]
\centering
\includegraphics[width=0.62\textwidth]{lecture_04_figure_05.png}
\caption{Two-gluon state count written as \(8\times 8=1+63\).}
\end{figure}

The blackboard first records the counting statement
\begin{equation}
8\times8 = 1+63.
\end{equation}
This is not yet a full irreducible decomposition. It is a checkpoint: two gluons give \(64\) states, one combination is a singlet, and the remaining \(63\) belong to non-singlet sectors.

Later in the lecture and discussion this is cleaned up into the more detailed representation-theoretic statement
\begin{equation}
8\otimes8 = 1\oplus8\oplus8\oplus10\oplus\overline{10}\oplus27.
\end{equation}
We should state this cautiously, because the board itself shows only the state count. But the important physical conclusion is clear enough: there is a singlet two-gluon state, and that is the lecture's glueball.

From here the lecture draws a decisive line. The free particles of the strong interactions are color singlets. In the language of the representation theory, only the \(1\) appears in the spectrum of isolated states. Baryons, mesons, and glueballs are the basic colorless hadrons.

A further payoff is electric charge. Because the free hadrons are built only from \(q\bar q\), \(qqq\), or colorless gluonic combinations, they carry only integer electric charge. That is one of the lecture's major motivations for color singlets: they explain why fractionally charged quarks do not appear as free particles.

\section{Why Only Singlets Survive}

At this point the lecture turns from classification to dynamics. This is where Nambu enters the story. His guiding idea was simple: if the low-energy states are forced to be color singlets, then fractionally charged particles never appear as liberated objects. But that leaves the real question: why should the dynamics of the theory favor singlets so strongly?

The first step is the analogy with electrodynamics. In a gauge theory, conserved quantities source photon-like fields. In electrodynamics the source is electric charge and the field is the electromagnetic field. In quantum chromodynamics the corresponding sources are the eight color generators, and the corresponding fields are the eight gluon fields. The lecture is careful here: color means the way an object transforms under the group, but the actual radiated field is sourced by the corresponding color charge.

The second step is the crucial difference between photons and gluons. Photons are neutral and therefore do not source one another. Classical electrodynamics is linear in that sense. Gluons, by contrast, carry color. The gluon field therefore interacts with itself.

This self-interaction is the lecture's dynamical turning point. The lines of color flux do not spread freely as electric field lines do in an ordinary dipole pattern. They attract one another and bundle into a tube of flux. The resulting field has approximately uniform energy per unit length, so the energy of a separated quark--anti-quark pair grows linearly with the separation:
\begin{equation}
V(r)\propto r.
\end{equation}
This is the lecture-level statement of confinement. It is not presented as a rigorous derivation with a calculated string tension, but as the central physical effect of the nonlinear gluon dynamics.

The consequences are immediate. A non-singlet object drags a color flux tube behind it. If we try to pull a quark away from an anti-quark inside a meson, the field energy keeps growing. Either the quark is pulled back, or the tube breaks by pair creation:
\begin{equation}
q\bar q \longrightarrow q\bar q + q\bar q.
\end{equation}
The system never produces an isolated quark. Instead it produces additional mesons. In the high-energy version of the story, this repeated breaking leads to a spray of hadrons, a jet.

The lecture also makes room for baryons in this picture. Three quarks may combine into a colorless object, and the flux lines can join so that there is no net color field extending to infinity. This is why the three-quark singlet is not just algebraically allowed but dynamically natural.

Finally, the lecture widens the frame again. Quantum chromodynamics is the \(SU(3)\) part of the Standard Model. Together with electrodynamics, which provides the \(U(1)\) comparison case, we have covered nine of the twelve conserved charges and gauge fields in the Standard Model: eight gluons and the photon. What remains is the additional \(SU(2)\) sector, along with the Higgs field, to be taken up in the next lectures.

\section{Summary}

The lecture begins with a question about groups and representations, and ends with a picture of confinement. In between, it builds a clean spine. First, the same abstract symmetry may act in different-dimensional representations depending on the system. Second, for \(SU(2)\) the infinitesimal generators are traceless Hermitian matrices, and in the rotation problem they are angular momentum operators. Third, the same machinery extends to \(SU(3)\), where quarks live in the \(3\), anti-quarks in the \(\bar 3\), and gluons in the \(8\). Fourth, the tensor products \(3\otimes3\), \(3\otimes\bar 3\), and \(3\otimes3\otimes3\) explain how color singlets arise as mesons and baryons, while \(8\otimes8\) allows glueballs.

Only then does the lecture ask the dynamical question: why do only singlets survive as free particles? The answer is the nonlinear dynamics of the gluon field, which squeezes color flux into tubes and makes isolated color prohibitively expensive in energy. In that sense the lecture gives us both the algebra of QCD and its first serious physical payoff.